\chapter{الگوریتم امضای دیجیتال \en{Ed25519}}
\label{app:ed25519}
\section{منحنی و تولید کلید}
\en{Ed25519} گونه‌ای مشخص از \en{EdDSA} است که بر منحنی ادواردز روی میدان با مشخصۀ $p=2^{255}-19$ تعریف می‌شود \cite{ed25519}. معادلۀ منحنی عبارت است از
\begin{equation}
-x^2+y^2=1+dx^2y^2,\qquad
d=-\frac{121665}{121666}\pmod p.
\end{equation}
نقطۀ پایۀ $B$ زیرگروه با مرتبۀ اول
\begin{equation}
\ell=2^{252}+27742317777372353535851937790883648493
\end{equation}
را تولید می‌کند و هم‌عامل منحنی ۸ است. این مرتبه با مشخصۀ میدان تفاوت دارد؛ نماها در پیمانۀ $\ell$ و مختصات در میدان $\F_p$ بررسی می‌شوند.

از یک بذر خصوصی ۳۲بایتی $z$، خروجی $h=\operatorname{SHA512}(z)$ محاسبه می‌شود. بخش نخست با عملیات تعیین‌شدۀ \en{clamping} به اسکالر خصوصی $a$ و بخش دوم به پیشوند محرمانۀ $h_{\mathrm{prefix}}$ تبدیل می‌شود. کلید عمومی، کدگذاری نقطۀ $A=[a]B$ است. جزئیات ترتیب بایت، کدگذاری نقطه و قیود بیت‌ها بخشی از استاندارد هستند و نباید با پیاده‌سازی دستیِ تقریباً مشابه جایگزین شوند.

\section{تولید امضا}
برای پیام بایتی $M$، اسکالر موقت و نقطۀ امضا چنین محاسبه می‌شوند:
\begin{align}
r&=\operatorname{SHA512}(h_{\mathrm{prefix}}\Vert M)\bmod\ell,\\
R&=[r]B,\\
k&=\operatorname{SHA512}(\operatorname{enc}(R)\Vert\operatorname{enc}(A)\Vert M)\bmod\ell,\\
s&=(r+ka)\bmod\ell.
\end{align}
امضا شامل کدگذاری $R$ و $s$ است و مجموعاً ۶۴ بایت دارد. قطعی بودن این ساختار، نیاز به حفاظت از بذر و پیشوند محرمانه را حذف نمی‌کند. استفادۀ نادرست از همان مواد کلید در گونه‌ای ناسازگار می‌تواند پیش‌فرض‌های امنیتی را تغییر دهد.

\section{بررسی و رابطۀ درستی}
در توضیح جبریِ زیرگروه معتبر، بررسی امضا به رابطۀ
\begin{equation}
[s]B=R+[k]A
\end{equation}
متکی است، زیرا
\begin{equation}
R+[k]A=[r]B+[ka]B=[r+ka]B=[s]B.
\end{equation}
در بررسی بایتی واقعی، صحت کدگذاری، محدودۀ اسکالر و ملاحظات هم‌عامل نیز اهمیت دارند. صورت‌های استانداردِ بررسی می‌توانند رابطۀ ضرب‌شده در هم‌عامل را به کار ببرند؛ این جزئیات در استاندارد و کتابخانۀ معتبر تعریف می‌شوند. رابطۀ جبری ساده، مجوز پذیرفتن هر رشتۀ بایتی به‌عنوان نقطۀ معتبر نیست.

\section{کاربرد در پروتکل پایان‌نامه}
در مخزن، کتابخانۀ \en{cryptography} عملیات \en{Ed25519} را انجام می‌دهد. پیامِ ورودی کتابخانه همان خروجی قطعیِ \en{canonical\_message} است. این پیام شامل شناسۀ طرح، کلاینت، دور، \en{nonce}، متن رمز و وزن است؛ امضای مستقل هر فیلد یا امضای یک هش دستیِ متفاوت انجام نمی‌شود.

این لایه از دستکاری برون‌زاد محموله و جعل هویت در چارچوب فرض‌های کلید جلوگیری می‌کند، اما محرمانگی متن را فراهم نمی‌کند. در مقابل، رمزگذاری بنالوح نیز به‌تنهایی هویت فرستنده را ثابت نمی‌کند. ترکیب درست نقش‌ها، بدون ضرب امضا در متن رمز، علت اصلی تفکیک دو مؤلفه در فصل چهارم است.
